\documentclass[11pt]{article}

\usepackage[margin=1.1in]{geometry}
\usepackage{amsmath,amssymb,amsthm}
\usepackage{booktabs}
\usepackage{microtype}
\usepackage{xcolor}
\usepackage[colorlinks=true,linkcolor=blue!60!black,citecolor=blue!60!black,urlcolor=blue!60!black]{hyperref}

\newtheorem{theorem}{Theorem}[section]
\newtheorem{proposition}[theorem]{Proposition}
\newtheorem{corollary}[theorem]{Corollary}
\newtheorem{lemma}[theorem]{Lemma}
\theoremstyle{definition}
\newtheorem{claim}[theorem]{Claim}
\newtheorem{openproblem}[theorem]{Open Problem}
\theoremstyle{remark}
\newtheorem{remark}[theorem]{Remark}

\newcommand{\PT}{\mathcal{PT}}
\newcommand{\dd}{\mathrm{d}}
\newcommand{\ii}{\mathrm{i}}
\newcommand{\ee}{\mathrm{e}}
\newcommand{\norm}[1]{\left\lVert #1 \right\rVert}
\newcommand{\avg}[1]{\langle #1 \rangle}
\newcommand{\ket}[1]{|#1\rangle}

\title{Ghost vacuum decay as a stochastic process:\\
instantaneous states, the boost-orbit measure,\\
and the foliation that ghosts make visible}

\author{Jonathan Torkelson\\[2pt]
\small Embeddetech, Inc., Tulsa, Oklahoma, USA\\
\small\texttt{jtorkelson@embeddetech.com}}
\date{July 11, 2026}

\begin{document}
\maketitle

\begin{abstract}
In a companion paper we showed that the Ostrogradski instability of
interacting ghosts is, in the stochastic currency of Barandes'
quantum--stochastic correspondence, a recurrence-classification
problem, with malicious interactions realizing \emph{explosion}:
escape from the sample space in finite time. Here we take the program
to quantum field theory, where the vacuum of a ghost coupled to
ordinary matter famously decays at a divergent rate. Our results:
(i)~the explosion dictionary acquires a Fock-space column---Feller's
\emph{pure-birth} ladder criterion $\sum_n 1/\lambda_n<\infty$ is the discrete image of
the classical time-of-flight integral, and finite-time blowup
corresponds to ladder exponents $\alpha_{\rm eff}>1$, verified numerically (up to
marginal logarithmic cases at $\alpha_{\rm eff}=1$); (ii)~the known factorization of the ghost
vacuum rate into an infinite boost-orbit volume times an invariant
reduced integral (Gross--Strumia--Teresi--Zirilli) is verified by
brute-force phase-space integration to $0.1\%$ in $d{=}1$, and yields
the reduced measure's dimension-and-coupling ladder
($\Lambda^{0}/\Lambda^{4}/\Lambda^{8}$),
whose $\Lambda^{8}/M_{\rm Pl}^{4}$ rung reproduces the
Cline--Jeon--Moore scaling; (iii)~in Markov language
the Lorentz-invariant \emph{eternal} ghost vacuum is, on the
multiplicity reading and at tree level, an
\emph{instantaneous state}---the minimal semigroup degenerates to
zero vacuum lifetime---so a coupled ghost admits a non-degenerate
configuration-space process only relative to a structure that individuates the channels---in every known rescue a frame or a birth surface---and every
configuration-preserving rescue in the literature (momentum cutoffs, ghost condensation, a
born vacuum) supplies one; (iv)~classical lattice experiments in
$1{+}1$ and $2{+}1$ dimensions show the classical decay tracks a
different, \emph{tadpole} dimension ladder ($\avg{\phi^2}\sim
k_{\max}^{d-1}$, logarithmic at $d{=}1$), with the channel ladder flat to classical dynamics within the
measurement resolution---spontaneous emission into empty channels is intrinsically
quantum; (v)~a beables-first field theory is foliation-anchored by
construction, the foliation's determinacy is a measurable property of
the state that vanishes continuously at the vacuum, and a ghost
sector would render the otherwise-hidden foliation observable (the hiddenness
proved for the finite core; the ghost-induced observability argued at
tree level);
(vi)~apparent superluminal signaling enters at the two points we
identify (the conditioning point provably the unique marginal-breaker in the finite core, over the framework's dilated local dynamics; the completions door left open), both simulated: survival conditioning (recovering the
known resolution of the Lee--Hsieh--Flammia--Lee no-signaling
controversy in one bookkeeping, with the signal magnitude of its
broken-phase dimer analog now in
closed form, $\mathrm{TV}(\infty)=\tilde\kappa/\sin\theta$)
and conservative completions (a messaging protocol whose latency
through infinity is distance-independent in the explosive
class (in the cases tested), and which the minimal process silences). The findings assemble
into a trilemma: \emph{coupled ghosts, empirical Lorentz
invariance, definite configurations---any two, never all three} (at
tree level).
\end{abstract}

\newpage
\begin{quote}
\small\noindent\textbf{Note on provenance.} This manuscript is an
experiment in AI-conducted research as much as a physics paper. The
project began when the author asked the large language model Claude
Fable~5 (Anthropic): \emph{``Is it not possible that the antilinear
PT/time-reversal symmetry element could maybe predict indivisible
non-Markovian processes?''} Everything that followed was performed by
the model: the research program, the theorems and proofs, the numerical
experiments, the literature search, and this text. The model was
repeatedly prompted to attack its own results adversarially and to
search for prior work. The author, a professional engineer rather than
a physicist, directed and reviewed the sessions but makes no claim to
comprehension of the findings, and is circulating the manuscript for
independent review by professional physicists. Accordingly, nothing
here should be taken on trust: every theorem is stated with a full
proof or its explicit verification, and every numerical claim is
reproducible from the public scripts (see Code availability); the
originating conversation, project briefing, and full research
notebook are preserved in the program's hub repository (see Code
availability).
\end{quote}

\newpage
\tableofcontents
\newpage

%=====================================================================
\section{Introduction}
\label{sec:intro}

The companion paper \cite{indivisible} developed a dictionary between
non-Hermitian and higher-derivative quantum mechanics and the
indivisible-stochastic-process reconstruction of quantum theory due
to Barandes \cite{Barandes2023,Barandes2025}. Its final result, the
\emph{explosion dictionary}, identified the Ostrogradski instability
of malicious interacting ghosts \cite{PaisUhlenbeck1950,Smilga2005}
with \emph{explosion} in the probabilist's sense: the classical time
of flight to infinity, the quantum deficiency norm, and the failure
of Feller conservativity are one integral, and the canonical
description of a malicious ghost is a \emph{sub-stochastic} process
that loses mass to infinity in finite time.

That paper stopped at quantum mechanics. The present paper asks
whether the identification survives the passage to quantum field
theory, where the instability of a ghost coupled to ordinary matter
takes its sharpest known form: the vacuum itself decays, into ghost
pairs plus matter pairs of balancing energy, at a rate whose
phase-space integral diverges \cite{CHT2003,CJM2004}. Whether that
divergence means catastrophe or artifact is a live dispute:
Cline--Jeon--Moore concluded that only a Lorentz-\emph{violating}
momentum cutoff tames it and bounded the cutoff at
$\Lambda\lesssim3\,$MeV from the diffuse gamma-ray background
\cite{CJM2004} (a 2023 re-analysis with exact phase space relaxes
this to $\Lambda\lesssim19\,$MeV \cite{ClinePhantom2024}); Gross--Strumia--Teresi--Zirilli, noting that the same
divergent boost integral once afflicted ordinary false-vacuum decay
before Coleman's bounce resolved it as an over-counting, answer
``we do not know the true rate'' \cite{GSTZ2021};
Garriga--Vilenkin construct Lorentz-invariant ghosts whose finite
vacuum age and UV-soft nonlocal couplings yield a finite,
time-growing rate \cite{GV2013}. The dispute matters beyond theory:
dark-energy measurements with $w<-1$---currently hinted at---would
make the phantom (ghost) sector a live physical hypothesis
\cite{CHT2003,ClinePhantom2024}.

This paper translates the dispute into the stochastic currency and
finds that the translation does real work. Our main results:

\begin{enumerate}
\item \textbf{The Fock-space column of the explosion dictionary}
(Sec.~\ref{sec:pathologies}). Feller's ladder criterion---a
\emph{pure-birth} jump process explodes iff $\sum_n 1/\lambda_n<\infty$---is
the discrete image of the companion paper's time-of-flight integral,
and the transient/explosive dichotomy is a statement about growth
\emph{laws}: algebraic blowup $n\sim(t_*-t)^{-m}$ has ladder exponent
$\alpha_{\rm eff}=1+1/m>1$; exponential growth sits exactly at the
marginal $\alpha_{\rm eff}=1$. Both verified numerically, including
the box-independence signature of explosion.
\item \textbf{Two pathologies, kept distinct}
(Sec.~\ref{sec:pathologies}). Explosion (finite rates growing along
the occupation ladder; escape at finite $t_*>0$) and
\emph{instantaneity} (infinite exit rate at a single state; zero
lifetime) are different diseases. Malicious quantum-mechanical ghosts
realize the first; the Lorentz-invariant ghost vacuum threatens the
second.
\item \textbf{The factorization, verified; the dimension ladder}
(Sec.~\ref{sec:factorization}). The vacuum rate factorizes as
[boost-orbit volume, infinite in every dimension] $\times$ [invariant
reduced integral over the pair invariant mass]---a structure due to
GSTZ \cite{GSTZ2021} (implicit in CJM's regulated calculation
\cite{CJM2004}), which we verify by brute-force constrained
integration to $0.1\%$ in $d=1$ and whose reduced measure we find numerically to be
UV-finite in $d=1$, $\sim\Lambda^{4}$ in $d=3$, and
$\sim\Lambda^{8}/M_{\rm Pl}^{4}$ with the CJM gravitational
vertex---their published scaling, recovered from the orbit-quotiented
measure alone.
\item \textbf{Instantaneity and the frame}
(Sec.~\ref{sec:multiplicity}). If boost-orbit members are distinct
outcomes with tree-level (orbit-degenerate) rates, the eternal
Lorentz-invariant ghost vacuum has infinite exit rate per unit
volume: an \emph{instantaneous state}, whose minimal semigroup
degenerates to zero vacuum lifetime. A born vacuum (finite age), with
a UV-finite \emph{causally weighted} reduced integral (a sufficiently
UV-soft coupling; in $d=1$ a logarithmically weighted finiteness,
Thm~\ref{thm:instantaneity}), has finite, growing rate \cite{GV2013}. Every known
rescue that preserves definite configurations supplies a frame or a birth
surface; in this currency, \emph{a
coupled ghost admits a non-degenerate sample space only relative to an
individuating structure---a frame or a birth surface}. We port the bubble literature's multiplicity-vs-redundancy
resolution \cite{KOV1975,ColemanBounce,Dvali2011,DDP2012,GKT2013,GSV2011}
and show its two mechanisms (saddle symmetry; excitation cost) are
absent, at tree level, for perturbative ghost decay.
\item \textbf{The two ladders} (Sec.~\ref{sec:lattice}). Classical
lattice experiments at fixed volume show the classical decay rate is
effectively marginal in $1{+}1$d (consistent with the logarithmic tadpole expectation, not cleanly separable from saturation at accessible $N$) and power-law divergent
in $2{+}1$d---but exponent
decomposition attributes the mechanism to the \emph{tadpole} ladder
($\avg{\phi^2}_0\sim k_{\max}^{d-1}$), not the channel ladder: with
the source held fixed, the channel factor is flat within resolution. Classical fields
cannot spontaneously radiate into empty channels; the channel ladder
is intrinsically quantum.
\item \textbf{The foliation} (Sec.~\ref{sec:foliation}). A
beables-first field theory chooses configurations on slices: it is
foliation-anchored by construction, as in Bohmian field theory
\cite{HypersurfaceBD,DGNSZ2014}. We show the foliation's
\emph{determinacy} is a measurable property of the state, vanishing
continuously at the vacuum (the dynamical face of
$\avg{T^{\mu\nu}}\propto\eta^{\mu\nu}$); that a ghost sector would
make the otherwise-hidden foliation \emph{observable}; and that the ultraviolet disposes in a
lattice experiment while the state proposes by continuum covariance (the
lattice could not resolve the early window at accessible sizes)---the
dynamical content of Garriga--Vilenkin's UV-softness condition.
\item \textbf{Causality and signaling} (Sec.~\ref{sec:causality}).
Apparent superluminal signaling enters this framework at the two
points we identify (the conditioning point provably the unique marginal-breaker in the finite core, over the framework's dilated local dynamics; the completions door left open), both simulated. Survival conditioning: the
Lee--Hsieh--Flammia--Lee signal \cite{LHFL2014} is reproduced,
located, and---in its broken-phase dimer analog---given a closed-form
magnitude, $\mathrm{TV}(\infty)=\tilde\kappa/\sin\theta$, as a corollary
of the companion paper's terminal-indivisibility theorem; the correctly
bookkept (dilated) dynamics is no-signaling to $10^{-16}$.
Conservative completions: a messaging protocol through infinity
whose latency is distance-independent in the explosive class (in the
cases tested) and which the minimal process silences---\emph{causality selects the
minimal process}.
\item \emergencystretch=2em\textbf{The trilemma} (Sec.~\ref{sec:trilemma}):
coupled ghosts, empirical Lorentz invariance, definite
configurations---any two, never all three (at tree level).
\end{enumerate}

Section~\ref{sec:tests} collects observational and experimental
tests; Sec.~\ref{sec:discussion} states the limits of every claim.
All numerics are reproducible from standalone scripts
(\nameref{sec:code}).

%=====================================================================
\section{Two pathologies of a stochastic process}
\label{sec:pathologies}

We use the correspondence in the form of \cite{Barandes2023,Barandes2025,indivisible}: configurations carry the ontology, and the
dynamics is a (generally indivisible) family of stochastic maps; for
open or unbounded systems the regulator-independent object is the
minimal, generally sub-stochastic, semigroup \cite{Feller1952}.

\subsection{The Fock-space column of the explosion dictionary}

Let occupation configurations form a ladder with total exit rates
$\lambda_n$ from occupation $n$. Feller's criterion states that a
\emph{pure-birth} chain \emph{explodes}---performs infinitely many jumps in
finite time, with mass leaving every finite region---iff the expected
total travel time to infinity is finite:
\begin{equation}
\sum_n \frac{1}{\lambda_n} < \infty .
\label{eq:feller}
\end{equation}
This is the exact discrete sibling of the classical time-of-flight
integral $T=\int^\infty\dd x/\sqrt{2(E-V)}$ of the companion
dictionary: both are sums of sojourn times to infinity. The bridge
to dynamics is a one-line lemma:

\begin{lemma}[growth law decides]
\label{lem:growthlaw}
For any pure-birth ladder (rates $\lambda_n>0$, mean-field
$\dot n=\lambda_n$) with growth $n(t)$, define
$\alpha_{\rm eff}=\dd\log\dot n/\dd\log n$. Algebraic finite-time blowup $n\sim(t_*-t)^{-m}$ gives
$\alpha_{\rm eff}=1+1/m>1$; exponential growth gives
$\alpha_{\rm eff}=1$. By \eqref{eq:feller}, for power-law ladders
$\lambda_n=n^\alpha$ ($n\ge1$) the explosive class is
$\alpha_{\rm eff}>1$ (logarithmic
refinements can shift the marginal case): finite-time blowup, not
growth rate, is what explosion detects.
\end{lemma}

Numerically (\texttt{qft\_bridge\_stage2.py}): truncated-ladder
master equations with $\lambda_n=n^\alpha$ show the explosion
signature of \cite{indivisible}---absorbed mass at fixed time
\emph{independent of the truncation}---switching on at $\alpha>1$
(absorbed $0.941$--$0.988$, nearly flat (within $\sim2\%$) across truncations $200$--$1600$,
for $\alpha=1.25,1.5$; decaying to zero for $\alpha\le1$); and the
lemma's exponents are measured as $1.0000$ for the quadratic ghost
cascade ($n=\sinh^2 gt$: forever-growing yet \emph{transient}) and
$1.172$ vs.\ the predicted $7/6$ for the classical quartic
Pais--Uhlenbeck blowup (energy proxy $n\sim(t_*-t)^{-6}$).

\subsection{Explosion versus instantaneity}
\label{sec:taxonomy}

Both pathologies are failures of conservativity, but they are
distinct and must not be conflated:

\begin{itemize}
\item \emph{Explosion}: every exit rate finite, but growing along
the ladder so that \eqref{eq:feller} converges; jumps accumulate and
mass exits at some finite $t_*>0$. This is a \emph{tail} phenomenon
in configuration space. Malicious quantum-mechanical ghosts realize
it \cite{indivisible}.
\item \emph{Instantaneity}: the exit rate \emph{at a single state}
diverges, $q=\sum_{\rm channels}(\text{finite rates})=\infty$. States
with $q=\infty$ are called instantaneous; chains containing them
exist (Kolmogorov's and Blackwell's examples
\cite{Blackwell1958,ChungBook}) but are not \emph{stable}: they
admit no right-continuous jump realization. The operationally sharp
statement is obtained from regulated families: if
$q(\Lambda)\to\infty$ as the regulator is removed, the minimal
semigroup degenerates to \emph{instantaneous killing}---the state's
survival probability vanishes for every $t>0$. (Zero lifetime is
harsher than the honest instantaneous-state chains, in which the
state is still occupied a positive fraction of time; the difference
is that the minimal ghost-vacuum process has no return path into
the vacuum.) Not nonexistence: a zero-lifetime state.
\end{itemize}

The mathematical neighbors of this section deserve naming: global
existence (non-explosion) of Bell-type jump processes for lattice
QFT was proved by Georgii and Tumulka under growth hypotheses on
the Hamiltonian \cite{GeorgiiTumulka}, and conservativity criteria
for minimal quantum dynamical semigroups by Chebotarev and Fagnola
\cite{ChebotarevFagnola}; a ghost sector violates exactly the
hypotheses of such theorems, and this paper's content is that the
violation \emph{is} the physics. The passage from quantum mechanics
to field theory, we now show, is the passage from the first
pathology to the threat of the second.

%=====================================================================
\section{The ghost vacuum rate: factorization and the dimension
ladder}
\label{sec:factorization}

Consider a ghost field coupled to ordinary matter, with vacuum decay
$\varnothing\to(\text{matter pair})+(\text{ghost pair})$: total
four-momentum zero, the ghost pair carrying the negative of the
matter pair's energy--momentum, so the common pair energy is
unbounded \cite{CHT2003,CJM2004}. Because the matter pair's
four-momentum $Q$ must equal the sign-flipped ghost pair's, the
golden-rule rate per volume organizes, with $\rho(s)$ the invariant
two-body phase-space density and $s=Q^2$,
\begin{equation}
\frac{\Gamma}{V}
=\Bigl[\underbrace{\textstyle\int \dd\mu_{\rm boost}}_{\infty}\Bigr]
\times
\Bigl[\underbrace{\textstyle\int \dd s\,
\sigma_d(s)\,\rho_n(s)\rho_g(s)\,|\mathcal M(s)|^{2}}_{\text{reduced
integral}}\Bigr],
\label{eq:factorization}
\end{equation}
where the first factor is the volume of the boost orbit of the pair
frame---\emph{infinite in every spatial dimension $d$}---and
$\sigma_d(s)=\tfrac12 s^{(d-1)/2}$ is the residual invariant measure,
the Jacobian in
$\dd^{\,d+1}Q=\tfrac12 s^{(d-1)/2}\,\dd s\;\sinh^{d-1}\!\eta\,\dd\eta\,
\dd\Omega_{d-1}$ (with $Q^0=\sqrt s\cosh\eta$,
$|\vec Q|=\sqrt s\sinh\eta$), the last two factors being the
unit-hyperboloid orbit measure $\dd\mu_{\rm boost}$ (the overall
$\tfrac12$ is a boost-normalization convention, stripped in the
reduced integral $I_s$ of Prop.~\ref{prop:ladder}). This factorization
is due to Gross--Strumia--Teresi--Zirilli (their Eq.~(88) and the
diagnosis ``the rate contains a $\dd K_0$ integral over the
non-compact Lorentz group'' \cite{GSTZ2021}); the regulated special
case is CJM's Eq.~(9) \cite{CJM2004}. Our contributions are a
numerical verification and the dimension ladder of the reduced
factor.

\begin{proposition}[factorization verified; dimension ladder]
\label{prop:ladder}
(\texttt{qft\_bridge\_stage3.py})
In $d=1$ with contact coupling and masses $(m_n,m_g)=(1,0.7)$,
brute-force integration of the constrained two-pair phase space with
rapidity cutoff $|\eta|<H$ gives $\Gamma(H)$ linear in $H$ with slope
equal to the independently computed reduced integral
($I_s=\int\dd s\,\rho_n\rho_g=0.6192$ to the script's finite-$s$ grid, converging to $0.6195$; slope $=I_s/(2\pi)^2$ in
the conventions used): ratios $0.999,1.000,1.001,1.000$ at
$H=1,2,3,4$.
The reduced integral's UV behavior, fitted to three decimals:
$\Lambda^{0}$ (UV-finite) in $d=1$; $\Lambda^{4}$ in $d=3$ with
contact coupling; $\Lambda^{8}/M_{\rm Pl}^{4}$ in $d=3$ with the CJM
gravitational vertex $|\mathcal M|^{2}\sim s^{2}/M_{\rm Pl}^{4}$---%
recovering the published CJM scaling from the orbit-quotiented
measure alone. Two closeouts
(\texttt{qft\_bridge\_closeout.py}): the $d=3$ factorization is
verified by direct Monte Carlo with the hyperboloid orbit measure
$\sinh^{2}\!\eta\,\dd\eta$ ($\Gamma/V_{\rm orb}$ constant within the
${\sim}3$--$7\%$ sampling errors across $H=0.5$--$1.5$; the $\Lambda^{4}$
rung measured directly, global-fit slope $4.3$ ($4.1$ over the upper interval) as threshold effects
fade); and the causal orbit cap $\gamma<t\sqrt s$ reproduces
Garriga--Vilenkin's $D=2$ logarithmic growth with
$(2\pi)^2\,\dd\Gamma/\dd\ln t=0.6193$ against $I_s=0.6192$ (the same
stripped convention as the slope above)---their result is
this factorization, regulated by the vacuum's age.
\end{proposition}

Two consequences frame everything that follows. First, the
$d$-dependence everyone debates lives entirely in the reduced factor;
the (unregulated) orbit volume diverges identically in every dimension. (A naive
scale-counting argument that ``$d=1$ converges'' samples a non-orbit
direction and misses a vanishing Jacobian along the boost flow; the
numerics adjudicated for the invariance argument. Our notebook
records the error.) Second, Garriga--Vilenkin's $D=2$ result---a
Lorentz-invariant rate, finite and logarithmically growing in the
vacuum's age \cite{GV2013}---is precisely the $d=1$ reduced
factor times a causally regulated (logarithmic) orbit.

%=====================================================================
\section{Multiplicity or redundancy: the boost orbit as a
sample-space question}
\label{sec:multiplicity}

Equation \eqref{eq:factorization} forces the question on which the
CJM/GSTZ dispute turns: do the members of a boost orbit count as
\emph{distinct decay channels} (rates add; the total diverges) or as
\emph{redundant descriptions} of one outcome (quotient; the rate is
the reduced integral)?

\subsection{The bubble precedent, and why it does not transfer}

The identical question has a fifty-year history in ordinary
false-vacuum decay. Kobzarev--Okun--Voloshin observed the divergent
boost integral \cite{KOV1975}; Coleman's Euclidean bounce resolved it
(the $O(4)$-symmetric saddle's boost directions are a stabilizer, not
a moduli space; integrating them would over-count
\cite{ColemanBounce}); Dvali revived the worry \cite{Dvali2011};
Dine--Draper--Park answered in Minkowski language: a boosted bubble is
a \emph{different, excited} configuration whose production amplitude
is suppressed, so the channel sum converges \cite{DDP2012}; and
Garriga--Shlaer--Vilenkin showed the fluctuation-dressed bubble state
is boost-invariant---boosts map it to itself \cite{GSV2011}. The
literature's operative criterion, due in substance to Coleman and
stated sharply by Garriga--Shlaer--Vilenkin \cite{GSV2011}: \emph{if
the final state is a boost singlet, the boost integral is
redundancy} (over-counting of one outcome). We use the complementary
direction---non-singlet finals are distinct channels---together with
the tree-level degeneracy of their amplitudes, which is what carries
the divergence; Dine--Draper--Park's bubbles show distinctness alone
does not suffice (a non-singlet channel sum can converge when the
amplitudes are suppressed \cite{DDP2012}). The born-vacuum idea was
itself anticipated in Dvali's note~5 \cite{Dvali2011}.

Perturbative ghost decay fails the criterion on both counts, at the
order computed: multiparticle final states with definite momenta are
not boost singlets, and orbit members are exactly degenerate in
$|\mathcal M|$---there is no saddle symmetry and no excitation cost.
(Whether resummation generates a Dine--Draper--Park-style suppression
is precisely GSTZ's open question restated; our statements are
tree-level.) On the multiplicity horn, then:

\begin{theorem}[instantaneity of the eternal ghost vacuum]
\label{thm:instantaneity}
If boost-orbit members are distinct outcomes \emph{with the
tree-level (orbit-degenerate) golden-rule rates} and a continuous, nonnegative reduced integrand not vanishing identically on the accessible range, the exit
rate \emph{per unit volume} of the Lorentz-invariant eternal vacuum of a coupled ghost
theory is infinite in every spatial dimension: the vacuum is an instantaneous
state (its total rate diverges already at any fixed finite volume), and the minimal semigroup of the regulated family degenerates
to zero vacuum lifetime as the regulator is removed. A vacuum
\emph{born} on a spacelike surface $\Sigma$ at finite age $t$, whose
\emph{causally weighted} reduced integral
$\int\dd s\,s^{(d-1)/2}\sigma_d(s)\,\rho_n(s)\rho_g(s)\,|\mathcal M(s)|^{2}$
(this causal-cap weight $s^{(d-1)/2}$, derived in the proof below for $d\ge2$, is a second factor of the same power as, but distinct in origin from, the measure $\sigma_d$; the $d=1$ case is logarithmic)
is finite---locally integrable (no on-shell pole in the accessible range) and ultraviolet-finite for a sufficiently UV-soft coupling (for $d=1$ the
weight is instead \emph{logarithmic}---the capped orbit factor is
$\operatorname{arccosh}(t\sqrt s)\sim\ln(2t)+\tfrac12\ln s$, so
finiteness is guaranteed by the log-weighted condition
$\int\dd s\,\sigma_1\rho_n\rho_g|\mathcal M|^{2}\,(1+|\ln s|)<\infty$
(both the plain reduced integral and its $\ln s$ weighting), strictly
stronger than plain UV-finiteness of the reduced integral; this
sufficient condition is not necessary once the accessible range reaches
$s\to0$, where the causal cap $s>1/t^2$ removes the $\ln s$ weight), has finite
exit rate per unit volume (growing as
$t^{\,d-1}$, logarithmically for $d=1$; e.g.\
$\Gamma/V\sim10^{-6}(\mu^{8}/M_{\rm Pl}^{4})(\mu t)^{2}$ for UV-soft
gravitational-strength coupling \cite{GV2013}); for a UV-local coupling
in $d\ge2$ the causal age bound regulates only the boost orbit, and the
reduced integral's own ultraviolet divergence (computed at $d=3$ in
Prop.~\ref{prop:ladder} as $\sim\Lambda^{4}$; the same plain counting gives
$\sim\Lambda$ at $d=2$, and $\sim\Lambda^{2}$ once the causal weight
$s^{(d-1)/2}$ controlling $q(t)$ is folded in) keeps $q=\infty$---any positive
power suffices.
\end{theorem}

\begin{proof}
Immediate from \eqref{eq:factorization}: the reduced integral is
positive whenever the reduced integrand does not vanish identically, and the orbit volume is
infinite; regulating the orbit (by a frame cutoff \cite{CJM2004} or
by the causal bound $\gamma<t\sqrt s$ \cite{CJM2004,GV2013}) gives
finite $q(\Lambda)$, or---when the \emph{causally weighted} reduced
integral is ultraviolet-finite---finite $q(t)$, with $q\to\infty$ as the
regulator is removed. (The causal cap $\cosh\eta<t\sqrt s$ leaves the
orbit integral $\int_0^{\eta_{\max}}\!\sinh^{d-1}\!\eta\,\dd\eta\simeq
(t\sqrt s)^{\,d-1}/(d-1)$ for $d\ge2$ and $\simeq\ln(t\sqrt s)$ for
$d=1$; the resulting factor $s^{(d-1)/2}$ migrates into the reduced
integrand, so it is the \emph{weighted} ultraviolet finiteness, not the
plain one, that controls $q(t)$ in $d\ge2$.) The minimal-semigroup statement is the standard limit of
survival probabilities $\ee^{-q t}\to0$.
\end{proof}

\begin{remark}[on the golden-rule inputs]
The reduced integrand uses the on-shell asymptotic densities
$\rho_n(s)\rho_g(s)$. Where the ghost admits no asymptotic particle
interpretation above threshold \cite{Buoninfante2026}, these are formal
spectral inputs rather than S-matrix data, and the conclusion is read
throughout as the vacuum's exit rate on the configuration sample space
(Section~\ref{sec:discussion}), not a decay width; the instantaneity
statement is about that configuration-space process, whose inputs
remain the on-shell densities.
\end{remark}

\begin{claim}[the sample space requires individuating structure]
\label{claim:frame}
On the multiplicity horn, a coupled ghost admits a non-degenerate
configuration-space stochastic process only relative to a structure
that individuates the channels: a frame (CJM's cutoff frame
\cite{CJM2004}; the ghost condensate's background
\cite{GhostCondensate}; Ho\v rava-type explicit breaking
\cite{Horava-ghosts}) or a birth surface (\cite{GV2013}, where
$\Sigma$ ``breaks Lorentz invariance spontaneously''). Conversely,
CJM's exclusion of Lorentz-invariant cutoffs is the statement that
the individuating data cannot be specified invariantly.
\end{claim}

The redundancy horn is GSTZ's hope: a Coleman-style boost-singlet
saddle for ghost decay, which they could not construct (no positive
Euclidean action) \cite{GSTZ2021}. In our language the two horns are
two \emph{sample spaces}: frame-anchored momentum configurations
versus boost-equivalence classes. That the dispute is a choice of
sample space---not a computational disagreement---is, we believe, its
sharpest available formulation, and Sec.~\ref{sec:foliation} shows
the beables framework itself takes a side.

\begin{remark}[scope: the Born--Markov concession]
\label{rem:bornmarkov}
The jump process is derived from the golden rule; the instantaneity
diagnosis is a \emph{translation} of the divergent perturbative rate
and inherits both its validity and its doubts. The framework's
contribution is that the translation exposes the disputed
assumption---channel individuation---as a statement about the sample
space, not that it independently settles the rate.
\end{remark}

A zero-dimensional channel model makes the dichotomy quantitative
(\texttt{qft\_bridge\_stage2.py}): a state coupled to $M$ identical
flat quasi-continua shows fitted decay rate tracking $M\Gamma_1$
within $3\%$ up to $M=32$ (unbounded growth: the multiplicity trend),
while a decaying channel measure $g_k\propto\ee^{-k/4}$ converges to
its predicted plateau ($0.0640$ vs.\ $0.0643$): channel counting is
the entire difference between ``catastrophic'' and ``metastable.''

%=====================================================================
\section{The two ladders: what classical field dynamics can and
cannot see}
\label{sec:lattice}

Does classical lattice field theory---the standard nonperturbative
tool---exhibit the channel divergence? We ran the direct experiment
(\texttt{qft\_bridge\_e4\_2d.py}): normal field $\phi$ and ghost
$\chi$ with coupling $\lambda\phi^{2}\chi^{2}$, correctly normalized
zero-point (``vacuum-mimic'') initial data, at \emph{fixed physical
volume} with the lattice cutoff $k_{\max}=\pi N/L$ swept; observable,
the time for the normal sector to absorb a fixed energy. With all
controls clean (the $\lambda=0$ runs silent; conserved-energy drift
well below thresholds (worst case ${\sim}3\times$ below at the top $d{=}2$ cutoff, ${\ge}10\times$ elsewhere)):
\begin{center}
\begin{tabular}{lcccc}
\toprule
& \multicolumn{4}{c}{rate $=1/t_{\rm thr}$ vs.\ lattice cutoff}\\
$d=1$ ($N=16,32,64,128$): & $0.013$ & $0.018$ & $0.024$ & $0.025$\\
$d=2$ ($N=16,32,48,64$): & $0.02$ & $0.1$ & $0.2$ & $0.4$\\
\midrule
\multicolumn{5}{l}{\footnotesize The two rows are measured at
\emph{different} configurations: $d=1$ at
$\lambda=0.25$, $L=32$,}\\
\multicolumn{5}{l}{\footnotesize $\Delta E=25$, $n=8$;
$d=2$ at $\lambda=0.1$, $L=16$, $\Delta E=10$, $n=60$. The claim is the
\emph{dichotomy},}\\
\multicolumn{5}{l}{\footnotesize not a cross-row rate comparison.}\\
\bottomrule
\end{tabular}
\end{center}
The $1{+}1$d rate is effectively marginal (its tadpole driver only logarithmic in $d{=}1$---slow, unbounded growth rather than a clean continuum limit; the rate $0.013\to0.025$ is still rising at $N{=}128$, though the growth over the final doubling is within the ensemble scatter at $n=8$, so the quoted uncertainties are ensemble estimates)---while
the $2{+}1$d rate grows as $\sim k_{\max}^{2.1}$ with no
saturation: a dimensional dichotomy in the pre-registered direction. (The
top-cutoff $d{=}2$ rate scatters ${\sim}20$--$40\%$ across vacuum realizations---the
$d{=}2$ rates are medians over $60$ (the $d{=}1$ row over $8$), quoted to one
significant figure ($d{=}1$ to two), with the threshold checked on a ${\approx}0.5$ grid; the
$d{=}1$ values likewise carry ${\sim}{\pm}30$--$50\%$ (upper-skewed) seed-to-seed scatter, and the fitted
exponent carries a ${\pm}0.2$ ensemble uncertainty.)

But exponent decomposition (\texttt{qft\_bridge\_decomp.py})
reattributes the mechanism. Holding the source fixed (zero-point
support truncated to $|k|\le K_C$, $N$-independent up to
$O((k\Delta x)^2)$ dispersion corrections (${\lesssim}3\%$)) and
sweeping only the available channels, the purified channel factor is
\emph{flat} in both dimensions (fitted exponents $-0.20$ ($d{=}2$) and
$+0.11$ ($d{=}1$); ensemble noise leaves room for a residual channel-log of at
most ${\sim}1.3\times$ across the sweep---so the naive golden-rule channel prediction (linear in the
cutoff, ${\sim}4\times$; a residual channel-\emph{log} would give the weaker
${\sim}2.2\times$) is excluded a fortiori); holding the channels fixed and scaling the source amplitude
$a$, the rate goes as $a^{4.96}$ ($d{=}2$) and $a^{4.45}$ ($d{=}1$) against the
perturbative $a^{4}$; and the reconstruction reproduces the growth to within ${\sim}30\%$ in the exponent: the measured
$\avg{\phi^{2}}_0\propto k_{\max}^{1.16}$ raised to \emph{half} the
$d{=}2$ amplitude power ($4.96/2=2.48$, since the rate scales as
$a^{4.96}$ while the variance scales as $a^{2}$) predicts
$k_{\max}^{2.9}$, over-predicting the observed ${\sim}2.1$: the sign and
rough size of the dimensional reversal are reproduced, the precise exponent
not.

\begin{claim}[the two ladders]
\label{claim:twoladders}
There are two dimension ladders. The \emph{tadpole ladder}: the
zero-point variance $\avg{\phi^{2}}=\int\dd^{d}k/(2\omega(2\pi)^d)
\sim k_{\max}^{d-1}$ (log at $d{=}1$, linear at $d{=}2$)---this is
what classical-statistical lattice dynamics tracks, and it is the
true mechanism of the marginal-versus-power-law dichotomy above. The
\emph{channel ladder}: the reduced decay measure of
Prop.~\ref{prop:ladder} plus the boost orbit---the CJM/GSTZ object.
Classical dynamics cannot see it, for the standard structural
reason: emission rates into a mode go as $(n+1)$ quantum
mechanically but as $n$ (stimulated only) in the classical wave
limit---the missing spontaneous ``$+1$'' is exactly the golden-rule
channel factor---so the classical cascade populates high $k$ only by
nonlinear local-in-$k$ flux, insensitive to the far cutoff. The
cutoff pathologies of classical-statistical simulations seeded with
zero-point noise are documented \cite{EGW2014}, including the
cutoff-dependent $\avg{\phi^2}$ renormalization in real-time
false-vacuum-decay lattices \cite{LatticeFVD2022}; our decomposition
quantifies them for the ghost problem.
\end{claim}

Methodological consequence, stated plainly: classical-statistical
vacuum simulations are tadpole-ladder probes; channel-ladder
questions require quantum treatment (our channel toy above, the
phase-space integrals of Prop.~\ref{prop:ladder}, or genuine quantum
simulation). Our notebook documents the two artifacts caught en route
(a vacuum-normalization error, retracted with scope; Verlet
integrator drift at marginal timestep firing the $\lambda=0$
control)---each alone would have
produced a publishable-looking false trend, and each was caught by an
inline control column.

%=====================================================================
\section{The foliation}
\label{sec:foliation}

\subsection{Beables anchor a slicing; ghosts make it visible}

A beables-first field theory needs configurations, and field
configurations live on \emph{spatial slices}: choosing the
configuration space is choosing a foliation. This is not peculiar to
the present program---it is the known structure of Bohmian field
theory \cite{HypersurfaceBD}, where the foliation may even be
extracted covariantly from the quantum state itself \cite{DGNSZ2014}
---and Barandes' own development stays on finite configuration
spaces, deferring exactly this point \cite{Barandes2025}. Within the
natural field-theoretic extension, then, the boost-orbit question of
Sec.~\ref{sec:multiplicity} is answered structurally: relative to the
theory's own foliation, orbit members are distinct
configurations---the multiplicity horn.

The standard objection to foliated ontologies is that the foliation
is idle: in quantum equilibrium it is empirically invisible
(no-signaling), pure metaphysical surplus. Our results replace
``invisible'' with a \emph{criterion}:

\begin{claim}[the equilibrium shield and its breaking]
\label{claim:shield}
\emph{In quantum equilibrium}, the foliation is empirically hidden
(the equilibrium shield): for Bohmian mechanics this
\emph{no-signaling} statement is Valentini's signal-locality theorem
\cite{Valentini1991}. Its finite-configuration core in the present
framework is Proposition~\ref{prop:shieldcore}, which assumes
\emph{neither} a bounded-below spectrum \emph{nor} an equilibrium
ensemble: marginal-independence holds for \emph{any} dilated local
unitary, with spectral positivity entering only in the interacting
extension (Open Problem~\ref{op:shield}; cf.\
Sec.~\ref{sec:causality}). One other equilibrium route to
foliation observability has been proposed: Bohmian arrival-time
statistics in EPRB configurations \cite{AmblardDrezet}---a
laboratory-scale mechanism resting on the contested measurability
of arrival times, disjoint from the ghost route argued here. A coupled ghost sector breaks the
shield: the vacuum-decay observables of
Sec.~\ref{sec:multiplicity}---the rate and the decay spectrum,
isotropic precisely in the individuating frame (with a cutoff edge in
that frame where a frame-cutoff rescue applies)---are
\emph{frame-dependent local observables} depending on the foliation
(the trilemma's operative sense, Claim~\ref{claim:trilemma}: some
observable depends on the foliation), not a superluminal signaling
channel. \emph{Ghosts are foliation detectors.} (The no-signaling
shield for the indivisible framework at full strength is a conjecture
with a defined proof target, Open Problem~\ref{op:shield}; the
breaking is Claim~\ref{claim:frame} plus kinematics.)
\end{claim}

\subsection{Determinacy: the foliation as a property of the state}

The vacuum is Lorentz-invariant: $\avg{T^{\mu\nu}}\propto
\eta^{\mu\nu}$ has no distinguished timelike eigenframe, so a
DGNSZ-style extraction \cite{DGNSZ2014} returns nothing. We upgraded
this algebra to a measured curve
(\texttt{s3\_foliation\_determinacy.py}): ensembles combining a
frame-symmetric zero-point part with boosted thermal occupations at
temperature $T$, frame estimator $\hat v=P/E$ per realization. At
$T=0$ the estimate is pure scatter (mean $0.001$, scatter $0.12$,
$\mathrm{SNR}\approx0$; at $T=0.01$ the mean even lands on the wrong
sign); the signal rises with the thermal share to
$\mathrm{SNR}\approx51$ at $T=3$. \emph{The foliation information
content of a state vanishes continuously at the vacuum}---which is
Theorem~\ref{thm:instantaneity} in ontological dress: the state that
determines no foliation is the state whose ghost dynamics is
ill-posed (one direction only---a frame is necessary, not sufficient, to
tame $q=\infty$). Conversely, any perfect-fluid-form state with $\rho+p\neq0$ determines a
frame---note the boundary, $\rho+p=0$, is the phantom divide
$w=-1$---so the actual cosmological state selects the comoving frame,
and CJM's pragmatic choice of the CMB frame for the cutoff
\cite{CJM2004} is what a state-extracted-foliation theory
\emph{prescribes}.

\subsection{The state proposes; the ultraviolet disposes}

Is the state-selected frame dynamically stable? A lattice experiment
(\texttt{foliation\_m3.py}): both sectors prepared as thermal baths
boosted to $v$ relative to the lattice (traveling-wave construction;
fixed spectral support, so the state is cutoff-independent), the
ghost bath hotter; the frame of the secular energy transfer measured
as $v^{*}=(\dd P/\dd t)/(\dd E/\dd t)$, which equals $v$ for
production isotropic in the state frame and $0$ for
lattice-anchored production. Result: a two-epoch structure. Deep in
the evolution the transfer is UV-directed (at $N{=}64$, $54\%$ of transferred
energy lands in the band $3<|k|\le6$, far above the bath support
$|k|\le1.2$, and ${\sim}83\%$ above the support altogether; at $N{=}128$--$256$
over $90\%$ lands above the support, much of it beyond $|k|{=}6$) and every band shows
$v^{*}\approx0$: the state's frame is erased en route to the
ultraviolet. In the early window, continuum covariance
\emph{guarantees} state-frame anchoring as a matter of principle
(the boosted evolution is the boost of the rest evolution, so
$v^{*}=v$ exactly in the continuum limit); the lattice experiment,
however, could not resolve it: small ensembles suggested a positive
early-window signal, but a $100$-run settlement showed no significant
positive signal ($\lesssim2\sigma$) at both $N{=}128$ and $N{=}256$ (screened means
$+0.04\pm0.05$ and $-0.03\pm0.10$; the slope-ratio estimator is
heavy-tailed, and runs with negligible early-window transfer must be
excluded or they dominate the raw mean). What the experiment establishes robustly is the
erasure side---the UV-directed cascade and the absence of the
state's frame in every band at late times---while providing no
positive lattice evidence of the covariant window at accessible
sizes.
\emph{The state proposes the foliation (by covariance); the
ultraviolet disposes (by measurement).}
This is the dynamical content of Garriga--Vilenkin's condition: their
finite-age frame can hold only because their couplings are UV-soft
\cite{GV2013}; in a UV-local theory the cascade hands the frame to
the regulator.

\subsection{The trilemma}
\label{sec:trilemma}

\begin{claim}[pick two]
\label{claim:trilemma}
Coupled ghosts; \emph{empirical} Lorentz invariance (no \emph{in-principle} observationally
accessible preferred foliation---``accessible'' meaning some observable depends on the foliation with nonzero effect, realized detectability being a separate, flux-dependent question; this is the operative sense, since a
beables ontology always carries a foliation, hidden or not, by
Sec.~\ref{sec:foliation}); definite configurations (beables): any
two, never all three. The born-vacuum construction \cite{GV2013}
sits on the first horn because its birth surface is
\emph{observable} through the frame-locked decay products, as are
the cutoff frame \cite{CJM2004}, the condensate background
\cite{GhostCondensate}, and explicit breaking \cite{Horava-ghosts}.
\begin{itemize}
\item coupled ghosts $+$ beables $\Rightarrow$ the foliation becomes
observable (Claim~\ref{claim:shield}); every configuration-preserving rescue on record
supplies and thereby exposes one;
\item empirical Lorentz $+$ beables $\Rightarrow$ no coupled ghosts
(healthy sectors; the ontological foliation exists but hides in
quantum equilibrium);
\item coupled ghosts $+$ empirical Lorentz $\Rightarrow$ no beables: the
definition-level rescues---$\PT$/no-ghost quantization
\cite{BenderMannheim2008} (which the companion paper showed exits the
sample space \cite{indivisible}) and Euclidean state-space surgery
\cite{HawkingHertog}.
\end{itemize}
The claim is a statement about combinations, verified against every
example we could find; it does not assert which pair nature chose. The incompatibility is a tree-level statement: were a resummed suppression to render the boost-orbit multiplicity sum convergent (a Dine--Draper--Park-style escape---GSTZ's open question, Sec.~\ref{sec:discussion}), the first horn would relax and a Lorentz-invariant coupled ghost with definite beables could survive.
Two of its three \emph{edges} are established prior art: ghosts demand a
preferred frame (\cite{CJM2004}, and theorem-level for stable NEC
violation \cite{DubovskyNEC}), and realist hidden-variable accounts of quantum
correlations demand one (Hardy's theorem \cite{Hardy1992} and its
descendants). The triangle---and the observation that the frame
demanded by ghost regularization and the foliation demanded by
beables can be \emph{the same object}, jointly observable---is, to
our knowledge, new.
\end{claim}

\begin{claim}[the divergence lives in the state]
\label{claim:state}
At tree level, the particle-number expectation per unit volume $\avg{N(t)}/V$ of the
evolving coupled-ghost vacuum carries the same
$[\text{orbit volume}]\times[\text{reduced integral}]$ factorization
as the exit rate \eqref{eq:factorization}: the number expectation
sums the same final-state phase space as the golden-rule rate and
boost-orbit members share a common particle number, so the orbit
multiplicity factors out exactly as there. With the invariant reduced
integral held ultraviolet-finite and nonvanishing (its ordinary UV divergence regulated
as usual), it diverges iff the multiplicity sum does---so the genuinely
instantaneous divergence is the boost-orbit one, an ontology-neutral
expectation invoking no beable configuration space, though it remains
conditional on the tree-level state. (Channel-toy verification,
\texttt{s5\_state\_divergence.py}: $\dd\avg N/\dd t$ tracks
$M\Gamma_1$ within $3\%$ to $M=32$ for the flat orbit measure,
versus a plateau for the decaying measure.)
\end{claim}

\begin{remark}[event ontologies inherit the dichotomy]
\label{rem:flash}
``Definite configurations'' above means configurations on slices;
the known foliation-free alternatives---event (``flash'') and
mass-density ontologies, as in relativistic GRW
\cite{Tumulka2006,Tumulka2020} (whose interacting extension is
delicate; cf.\ the no-go of \cite{JGB2019})---have rates that are
\emph{functionals of local quantum-state content}. By Claim~\ref{claim:state}, they do
not escape the trilemma; they import its central question: on the
multiplicity reading the state acquires unbounded particle content
and any content-proportional event rate diverges with
it---instantaneity in event clothing---while a covariant flash law
that takes the redundancy reading must dispute the tree-level state
itself, which is GSTZ's open question once more. The
global-existence (non-explosion) theorems for Bell-type QFT jump
processes \cite{GeorgiiTumulka} make the point citable: their
hypotheses on the Hamiltonian are exactly what a ghost sector
violates. The flank thus narrows from ``event ontologies might
escape'' to ``event ontologies face the same multiplicity-or-redundancy
fork,'' with the
resummation question open for all parties equally.
\end{remark}

%=====================================================================
\section{Causality and signaling}
\label{sec:causality}

Where, if anywhere, can superluminal signaling enter this framework?
We identify two entry points, and both were simulated.

\subsection{Conditioning: the Lee--Hsieh--Flammia--Lee signal,
located and quantified}

Ref.~\cite{LHFL2014} showed that local $\PT$ evolution on one arm of
an entangled pair, treated as fundamental, permits superluminal
signaling. The controversy has a resolution literature we build on:
the apparent signal was experimentally located in the post-selected
subspace \cite{Tang2016}, the metric-corrected treatment restores
consistency \cite{Brody2016}, and the no-go theorems survive proper
embedding \cite{Ju2019}---our treatments (iii), (i), and (ii)
respectively. What the stochastic currency adds is all three in one
classical-probability bookkeeping, plus the closed form below. In
this currency
(\texttt{s1\_no\_signaling.py}; Bell pair, Alice's message the loss
orientation of a dimer, signal the total-variation distance between
Bob's marginals): (i)~metric-corrected closed dynamics, max TV
$1.7\times10^{-16}$; (ii)~dilated open dynamics without
post-selection, max TV $4.4\times10^{-16}$---the framework's
falsifiable core, clean; (iii)~survival-conditioned dynamics: the
signal appears, with closed-form magnitude.

\begin{theorem}[the conditioning signal, in closed form]
\label{thm:lhfl}
For the Bell state and the broken-phase dimer of \cite{indivisible}
(loss rate $\sin\theta$, coupling $s$---the companion's dimer
notation, unrelated to the invariant mass $s$ of Secs.~3--4---%
$s<\sin\theta$, so $\tilde\kappa=\sqrt{\sin^{2}\theta-s^{2}}$ is real), with
$p=c+\sin\theta\,\sigma$, $q=c-\sin\theta\,\sigma$, $r=s\sigma$
($c=\cosh\tilde\kappa t$, $\sigma=\sinh(\tilde\kappa t)/\tilde\kappa$;
this dimer $\sigma$ is unrelated to the residual measure $\sigma_d(s)$,
and the amplitude $q$ here to the exit rate $q$ of the explosion
dictionary above),
the survival-conditioned signal between the two loss orientations is
\begin{equation}
\mathrm{TV}(t)=\frac{|p^{2}-q^{2}|}{p^{2}+q^{2}+2r^{2}},
\qquad
\mathrm{TV}(t\to\infty)=\frac{\tilde\kappa}{\sin\theta}
=\sqrt{1-\Bigl(\frac{s}{\sin\theta}\Bigr)^{2}} :
\label{eq:lhfl}
\end{equation}
zero at the exceptional point, approaching unity deep in the broken
phase.
\end{theorem}

\begin{proof}
$(C\otimes\mathbf 1)\ket{\psi_{\rm Bell}}$ gives Bob's conditioned
marginal $p_B(b)=\norm{C\ket b}^{2}/\sum_i\norm{C\ket i}^{2}$; the
squared column norms of the broken-phase contraction are
$\ee^{-2\sin\theta\,t}(p^{2}+r^{2})$ and
$\ee^{-2\sin\theta\,t}(q^{2}+r^{2})$ by the closed forms of
\cite{indivisible}; the two messages exchange $p\leftrightarrow q$.
The asymptote follows from $c,\tilde\kappa\sigma\to
\tfrac12\ee^{\tilde\kappa t}$ and
$\tilde\kappa^{2}+\sin^{2}\theta+s^{2}=2\sin^{2}\theta$.
Numerics: \eqref{eq:lhfl} matches to $2.2\times10^{-16}$; the
$s$-sweep matches $\tilde\kappa/\sin\theta$ to six digits away from
the exceptional point.
\end{proof}

(The unbroken-phase construction of \cite{LHFL2014} itself is
treatment (iv) in the script: conditioned signal up to $0.577$,
identically zero under treatment (i) with the same Hamiltonians.)
The interpretation is exact: the ``FTL'' of local $\PT$ lives only in
treating the survival-conditioned (sub-stochastic, renormalized)
process as the whole sample space---and even there it is not
operational, since constructing the conditioned ensemble requires
Alice-side survival information that travels classically. Same
Hamiltonians, three bookkeepings, signal in exactly one.

\subsection{Completions: messaging through infinity}

The companion paper showed conservative completions of an explosive
process return probability from infinity after the finite classical
round trip \cite{indivisible}. We upgraded the observation to a
communication protocol (\texttt{s2\_infinity\_channel.py}): Alice
gates the right exit of $V=-x^{4}$ (explosive) or $V=-x^{2}$
(marginal) with a barrier (her bit); Bob is a direction-resolved
current detector near the left boundary, triggering only on flux
entering from infinity; the completions are periodic identification
(conservative ``wormhole'') versus absorbing layers (minimal
process). All twenty-four runs (two potentials $\times$ two completions $\times$ three separations $L=8,10,12$ $\times$ Alice's two bit settings) are clean
(the $L=6$ geometry is degenerate---Alice's station overlaps the
launch packet---and is excluded, as recorded in the notebook):

\begin{center}
\begin{tabular}{lcc}
\toprule
 & conservative (wormhole) & minimal (absorbing)\\
\midrule
$V=-x^{4}$ (explosive) &
\begin{tabular}{@{}c@{}}one-bit channel;\\latency
$0.18/0.15/0.16$ at $L=8/10/12$:\\ \textbf{distance-independent}
\end{tabular} &
\begin{tabular}{@{}c@{}}silence
($J_{\max}\le3\times10^{-5}$);\\norm $\to0.000$\end{tabular}\\
\midrule
$V=-x^{2}$ (marginal) &
\begin{tabular}{@{}c@{}}channel exists; latency\\
$0.94/1.02/1.10$: \textbf{grows with }$L$\end{tabular} &
silence; \emph{in-box} norm $\to0.000$\\
\bottomrule
\end{tabular}
\end{center}

The marginal potential $V=-x^{2}$ is essentially self-adjoint
(limit-point at infinity), so its true minimal process is
\emph{conservative}; the ``$\text{norm}\to0.000$'' in that row is the
in-box (finite-window) norm as probability leaves the simulation box,
and its wormhole/absorbing variants are box-boundary \emph{controls},
not genuine self-adjoint completions. Only $V=-x^{4}$ is limit-circle,
with real self-adjoint-extension freedom---so the three-completion
taxonomy below is a statement about $-x^{4}$.

A third completion completes the taxonomy
(\texttt{s2b\_dirichlet.py}): hard Dirichlet walls fire the
\emph{right-side} direction-resolved detector (reflection off
infinity, $t_B\approx0.23$, distance-independent, left detector
silent, norm conserved)---so the three completions of one interior
Hamiltonian (the explosive $V=-x^{4}$) produce three distinct
signatures: return from the
left (wormhole), return from the right (walls), silence (minimal).
The same interior Hamiltonian thus contains a distance-independent
messaging channel through infinity, a distance-penalized one, or none
at all, according to the \emph{completion}---regulator data, not
physics in $H$. That extension choices are experimentally
distinguishable has precedent at the \emph{origin}-singularity
anomaly (graphene supercritical Coulomb ladders \cite{Ovdat2017});
the taxonomy here is its counterpart at infinity. The model is nonrelativistic---there is no light cone, and
``acausal'' here means latency that fails to scale with separation,
the model's stand-in for light-cone violation. In that operational
sense, the acausal channel occurs---in the cases tested ($-x^4$ and
$-x^2$), and as the time-of-flight dictionary predicts---exactly in the
explosive class (the dictionary's time-of-flight boundary, visible as a
flat versus rising latency curve): \emph{causality selects the minimal
process}.
For a malicious ghost, probability conservation and causality cannot
both be kept in the completions tested.

\subsection{Synthesis: positivity as the pillar}

Assembling Secs.~\ref{sec:foliation} and~\ref{sec:causality}: within
a beables ontology, apparent superluminal signaling enters at the two
points we identify (the conditioning point provably the unique
marginal-breaker in the finite core, over the framework's dilated local
dynamics (Prop.~\ref{prop:shieldcore}); the completions door and
field-theoretic exhaustiveness remain open)---conditioning
fallacies and acausal completions. The
acausal-completions door and the equilibrium shield are gated by
\emph{spectral positivity} (bounded-below sectors give conservative
processes and no channel; a ghost sector opens them at once); the
conditioning door, by contrast, is a post-selection artifact present
even in bounded-below lossy systems (Theorem~\ref{thm:lhfl}), gated by
whether one conditions on survival rather than by positivity. The
finite-configuration core of the shield---no-signaling for \emph{any}
dilated local dynamics, positivity entering only in the interacting
extension (Open Problem~\ref{op:shield})---is a theorem:

\begin{proposition}[no-signaling core of the shield]
\label{prop:shieldcore}
Let the joint configuration space be $C_A\times C_B$ (finite), let
Alice's operation be any dilated local dynamics given by a
\emph{unitary} $U_A$ on $A\otimes\mathrm{anc}_A$, Bob's any
\emph{unitary} $U_B$ on $B\otimes\mathrm{anc}_B$, and let the joint
process be $W=|U_A\otimes U_B|^{2}$ from an arbitrary joint amplitude
$\Psi$, extended by the fixed initial ancilla configurations, with
Bob's marginal summed over his ancilla configurations. Then Bob's
marginal at every time is independent of Alice's choice of $U_A$.
\end{proposition}

\begin{proof}
$p_B(b)=\sum_{a',\beta}|\langle a',b,\beta|(U_A\otimes U_B)\Psi\rangle|^{2}
=\langle\Psi|(U_A^{\dagger}U_A)\otimes\bigl(U_B^{\dagger}(|b\rangle\langle
b|\otimes\mathbf 1_{\mathrm{anc}_B})U_B\bigr)|\Psi\rangle
=\langle\Psi|\mathbf 1\otimes U_B^{\dagger}(|b\rangle\langle
b|\otimes\mathbf 1_{\mathrm{anc}_B})U_B|\Psi\rangle$, with $a'$ running over Alice's configurations
\emph{including} ancilla (leak) configurations and $\beta$ over Bob's ancilla. Survival
conditioning breaks exactly the first equality (a projector is
inserted on the $A$ side), which is where---and only
where---Theorem~\ref{thm:lhfl}'s signal lives.
\end{proof}

The closed-unitary case of Proposition~\ref{prop:shieldcore} is the
no-signaling reading of the causal-locality condition Barandes, Hasan,
and Kagan impose \cite{BHK2025} (whose marginal-factorization equations
give it directly); our statement extends it to dilated (open, lossy) local dynamics
and identifies survival conditioning as the unique breaking point.

What remains open is the field-theoretic and statistical extension:

\begin{openproblem}[the shield theorem, full strength]
\label{op:shield}
Extend Proposition~\ref{prop:shieldcore} to field configurations
with interacting joint dynamics and quantum-equilibrium ensembles:
prove that spectral positivity of all interacting sectors implies
no-signaling with respect to the foliation (invisibility of the
slicing), and that the implication is sharp---i.e., any coupled
sector unbounded below yields foliation-dependent observables (the
converse direction of Claim~\ref{claim:shield}).
\end{openproblem}

%=====================================================================
\section{Observational and experimental probes}
\label{sec:tests}

The trilemma is empirically loaded through its conjunctions.

\emph{(a) Dark energy.} Robust confirmation of $w<-1$ would make
the ghost hypothesis live and activate the rest of this section;
$w\ge-1$ retires it as the dark-energy explanation. Currently DESI DR2 prefers evolving dark
energy over $\Lambda$CDM at $2.8$--$4.2\sigma$ \cite{DESIDR2} (see
\cite{DESIreview} for a review)---with the honest caveat that the
quoted significance depends on which supernova sample is combined, a
principled combination giving ${\sim}3.1\sigma$ \cite{DESIskeptic}.

\emph{(b) The MeV gamma-ray background, with a frame signature.} On
the multiplicity horn, a coupled phantom fills the sky with photons
up to the cutoff, already bounded by the diffuse gamma-ray data
\cite{CJM2004,GV2013,ClinePhantom2024}; the current bound is
$\Lambda\lesssim19\,$MeV \cite{ClinePhantom2024}, which places the
spectral edge of a bound-saturating phantom inside the unobserved
``MeV gap.'' The stochastic reading adds
the \emph{foliation signature}: the background is isotropic in the
individuating frame, so a detector moving at
$\beta=v/c=1.23\times10^{-3}$ sees a Compton--Getting dipole
\cite{CG1935} aligned with the CMB dipole (the kinematic-dipole
toolkit for cosmological backgrounds is standard
\cite{KachelriessSerpico}, including threshold-dependent amplitudes
\cite{KinDipole2026}) \emph{and} a
direction-modulated spectral edge
$\Lambda(\theta)\simeq\Lambda(1+\beta\cos\theta)$. A mock-sky
pipeline (\texttt{s4c\_dipole\_template.py}; exact aberration and
Doppler transport, estimators verified against their analytic
amplitudes) gives the template: count-rate dipole amplitudes
$3.7$--$4.3\times10^{-3}$ depending on the threshold observable
(all-photon rate $3\beta$; thresholded counts add
$\beta E_{\rm thr}/(\Lambda-E_{\rm thr})$; each mock photon carries
the arrival-rate weight $1+\beta\cos\theta$---an unweighted mock
measures only the aberration part $2\beta$ and is blind to the rate
factor, an error caught in our own audit and documented in the
notebook), requiring
$\mathbf{4.0}$--$\mathbf{5.5\times10^{6}}$ \textbf{signal
photons for }$5\sigma$ with the apex fixed by the CMB dipole; and
the edge modulation ($0.12\%$)---whose cosmologically integrated
arriving spectrum is a linear \emph{wedge} rather than a sharp step
(\texttt{s4d\_cosi\_forecast.py}), costing ${\times}5$--$10$ in
counts---is recoverable \emph{through}
realistic $3\%$ instrument resolution by percentile regression at
$\sim3\sigma$ with a few$\,\times10^{6}$ photons ($5\sigma$ near
$10^{7}$). Dilution: for
signal fraction $f_s$ of collected diffuse counts, $N$ scales as
$1/f_s^{2}$---and CJM's bound is set precisely by the signal
\emph{saturating} the observed diffuse background, so \emph{at} the
bound $f_s=O(1)$ and the measurement is exposure-limited. This is a
knife-edge, not a generic expectation: $\Gamma\propto\Lambda^{8}$
means a cutoff at half the bound gives $f_s\sim4\times10^{-3}$ and
inflates $N$ by $\sim6\times10^{4}$. Two further practical caveats:
the measured GeV-band diffuse dipole is percent-level and dominated
by non-kinematic (astrophysical) components
\cite{DGBdipole}, so extracting a $10^{-3}$ kinematic dipole
requires exposure-uniformity control at that level; the edge
observable is less exposed to both issues. The discriminant against mundane comoving
backgrounds (which share the Compton--Getting dipole) is the edge
modulation together with the absence of source correlation.

\emph{(c) Short-range gravity.} Garriga and Vilenkin already note that
their cascade bound ($\mu\sim10^{-3}\,$eV) is only marginally consistent
with $0.1\,$mm tests of the inverse-square law \cite{GV2013,Kapner2007},
leaving a narrow surviving Lorentz-invariant window near
$\mu\approx2\times10^{-3}\,$eV; improved torsion-balance experiments can close it,
eliminating the born-vacuum rescue outright.

\emph{The decision matrix.} The decay signal is a \emph{joint}
observable: vacuum decay requires genuinely negative-energy states
(a modified-gravity mimic of $w<-1$ produces no such background),
and the signal's frame-locking \emph{is} the foliation. One positive
measurement therefore confirms both hypotheses at once, and the
residual degeneracies close with one added axis---structure-growth
and perturbation data distinguishing a genuine phantom field from a
gravitational mimic:

\begin{center}
\begin{tabular}{@{}p{0.3\linewidth}p{0.3\linewidth}p{0.32\linewidth}@{}}
\toprule
 & \textbf{signal seen} (frame-locked dipole $+$ edge) &
 \textbf{signal excluded} (window closed)\\
\midrule
\textbf{phantom field established} ($w<-1$ $+$ growth data) &
ghosts \emph{and} foliation both confirmed, in one measurement &
ghost real but no \emph{observable} decay: evidence \emph{against} definite
configurations \emph{if} the frame cutoff saturates the diffuse-background
bound---then the exit-the-sample-space horn
\cite{BenderMannheim2008,HawkingHertog} wins (at tree level)\\
\midrule
\textbf{phantom excluded} &
(inconsistent under the phantom-DE identification; a signal would otherwise indicate a non-DE ghost sector) &
the observable ghost hypothesis retired; foliation undecided\\
\bottomrule
\end{tabular}
\end{center}

Both axes of the matrix are simulation-backed. The rows
(\texttt{s4a\_books\_balance.py}): identical lattices, identical
temperatures, identical coupling strength, only the partner
sector's energy sign flipped---the healthy pair \emph{equilibrates}
(transfer saturates at $7.4\pm1.0$, the equipartition scale, for
the entire second half of the run), while the ghost pair runs away
without bound (past $10^{3}$ by $t\approx120$; no joint equilibrium
exists \cite{GSTZ2021}): only a genuine phantom feeds an unending
flux. The columns (\texttt{s4b\_horn3\_silence.py}): one two-mode
model, three quantizations---the ghost quantization reproduces the
exact decay law $\avg n=\sinh^{2}(gt)$ to five decimals through
$gt\lesssim1.2$ (truncation tail thereafter); the
positive-spectrum requantization gives \emph{machine-exact zero} at
every time (the vacuum is an eigenstate); the healthy control shows
only bounded virtual pairs. A decay signal exists in exactly one
quantization.

Three of the four cells are informative---one (phantom established,
signal seen) decisive for both questions, a second (phantom established,
signal excluded) decisive for both only where the predicted flux is
detectable (see below), the third (phantom excluded, signal excluded) for the ghost question
alone; the fourth is inconsistent under the phantom-DE identification. The signal-excluded verdict
\emph{against} beables is decisive only where the predicted flux is
detectable: a frame-cutoff phantom whose cutoff $\Lambda$ sits well
below the diffuse-background bound decays unobservably slowly with its
configurations intact, and $\Gamma\propto\Lambda^{8}$ makes this
\emph{undetectable}-$\Lambda$ region wide (the \emph{detectable} band
above being the knife-edge of the dilution estimate), leaving that cell
undecided for the foliation question.
Conditional on the top row---exactly the measurement DESI-era
surveys are making---the gamma-ray program is a \emph{two-sided}
test of the foliation: the surviving window is already pinched to
$\mu\approx2\times10^{-3}\,$eV between the cascade bound and
short-range gravity \cite{GV2013}, so ``window closed'' is a
realistic outcome, and it would count against the born-vacuum rescue
outright (the low-$\Lambda$ frame-cutoff region remaining undecided, as
above). The bottom-right cell states the honest
floor: in a ghost-free equilibrium universe the foliation is
undecidable \emph{in principle}, because the equilibrium shield is
quantum mechanics' own no-signaling property---an experiment
exposing a hidden foliation there would contradict quantum
predictions themselves. (One speculative shield-breaker in the
literature is primordial quantum \emph{non}-equilibrium
\cite{Valentini2007}, which would expose the foliation without
ghosts; the contested Bohmian arrival-time route
(Sec.~\ref{sec:foliation}) \cite{AmblardDrezet} is another. We note
both and set them aside.) Given ghosts, nature's choice
among the trilemma's horns is a measurement program, not a
philosophy seminar.

\emph{(d) Analog platforms.} The ladder criterion
(Lemma~\ref{lem:growthlaw}) is realizable in cavity/circuit QED:
two-mode squeezing already is the transient cascade
($\avg n=\sinh^{2}gt$ is routinely measured), and engineered
superlinearly Bose-enhanced couplings would exhibit the transition
toward finite-time runaway at $\alpha_{\rm eff}>1$. The
frame-competition experiment of Sec.~\ref{sec:foliation} (state frame
versus medium frame of an instability's emission) is native territory
for polariton fluids and analog-gravity systems. The companion
paper's photonic tests (the terminal-indivisibility closed forms in
lossy couplers) probe the quantum-mechanical layer of the same
program \cite{indivisible}.

%=====================================================================
\section{Discussion: what is established, and at what strength}
\label{sec:discussion}

\emph{Theorem-grade}: the Fock ladder criterion and growth-law lemma;
instantaneity on the multiplicity horn
(Thm.~\ref{thm:instantaneity}); the conditioning-signal closed form
(Thm.~\ref{thm:lhfl}); the no-signaling core
(Prop.~\ref{prop:shieldcore}). \emph{Measured}: the factorization verification and the dimension ladder (Prop.~\ref{prop:ladder}); the completions table; the
tadpole/channel decomposition; the two-epoch frame competition; the
determinacy curve. \emph{Argued}: the final-state-invariance
criterion's verdict on tree-level ghost decay; frame individuation;
the state-content factorization (Claim~\ref{claim:state});
the trilemma (a verified pattern over known examples).
\emph{Conjectured}: the shield theorem at full strength (Open
Problem~\ref{op:shield}; its finite core is proved).

Limits, stated without decoration. Everything golden-rule--based
inherits perturbation theory (Remark~\ref{rem:bornmarkov}); the
degeneracy of orbit members is a tree-level statement, and a
resummed suppression would render the multiplicity sum convergent
(a Dine--Draper--Park-style escape, distinct from redundancy)---that
is GSTZ's open question, not ours to close. The lattice
results are classical-statistical, with the stated blindness to
spontaneous emission; the factorization is numerically verified to
$0.1\%$ in $d=1$ and to ${\sim}3$--$7\%$ by Monte Carlo in $d=3$
(Prop.~\ref{prop:ladder}), still higher $d$ resting on the invariance
argument and GSTZ's analytics. The foliation claims live within the natural
field-theoretic extension of the beables framework, which its author
has not committed to; and no experiment distinguishes the ontology
itself---only its conjunction with ghosts is testable. Finally, the
trilemma does not say which two nature picked; it says the three
advertised virtues of fundamental physics cannot all be kept, at tree
level, if ghosts exist, and it prices each exit.

\emph{Adjacent results, positioned.} The active
ghost-stabilization program of Deffayet et al.\ proves benign
behavior for finite-dof quantum mechanics \cite{Deffayet2026} and
demonstrates it numerically for classical $1{+}1$d fields \cite{Deffayet2504}; neither engages the
boost orbit, and both are consistent with our taxonomy (the QM
constructions live in the companion's benign/transient classes; the
$1{+}1$d stability sits at the mild end matched by our marginal
(logarithmic) $d{=}1$ rate).
Buoninfante's analysis \cite{Buoninfante2026}---ghosts above
threshold admit no asymptotic particle interpretation---is the
sharpest operator-theoretic challenge to any ``ghost decay''
framing; our object is the vacuum's exit rate on a configuration
sample space rather than an S-matrix ghost width, but if the
dressed vacuum ultimately supports no decay interpretation, that
would favor the convergent readings and we flag it. K\"onnig et
al.'s slow heavy-ghost decay \cite{Konnig2016} and Kaplan--Sundrum's
energy-parity construction \cite{KaplanSundrum} are horn-one
examples (frame-anchored cutoffs); Babichev's identification of
ghost decay with Cherenkov radiation \cite{Babichev2024} makes the
frame-anchoring vivid; the modern $\Gamma/V$ formalism is
\cite{Andreassen2016}; and the information-flow criticality of
$\PT$ systems \cite{KAU2017} is the quantum-state-level cousin of
the conditioning analysis of Sec.~\ref{sec:causality}.

The companion paper ended by observing that the ghost problem, in the
stochastic currency, is a recurrence-classification problem. The
field-theoretic sequel is that it is, more deeply, a
\emph{sample-space} problem: the divergent vacuum rate, the
regulator dependence, the foliation, and every apparent causality
violation we examined are one question---\emph{which configurations
are distinct, and relative to what}---asked at increasing energy.
Every pathology entered through that door, and every resolution on
record pays at that door.

%=====================================================================
\phantomsection
\section*{Author contributions and use of AI}
\addcontentsline{toc}{section}{Author contributions and use of AI}

This manuscript was produced through a collaboration, directed by the
author, with the large language model Claude Fable~5 (Anthropic). The
AI system carried out essentially all of the technical work: the
research program and its strategy, the derivations and proofs, the
design, execution, and interpretation of all simulations, the
literature analysis, and the drafting of the manuscript, including
this statement. J.T.\ posed the research directions, provided
iterative review and quality pressure, and takes sole responsibility
for the decision to circulate the manuscript. In line with prevailing
editorial policies, the AI system is not listed as an author. The
research notebook in the public repository records the work
hour-by-hour, including three same-day corrections of our own errors
and one retraction, so that the process as well as the results can be
audited.

%=====================================================================
\phantomsection
\section*{Code availability}
\label{sec:code}
\addcontentsline{toc}{section}{Code availability}

This paper's scripts, source, and notebook excerpt are public at
\begin{center}
  \url{https://github.com/embeddetech/p03-instantaneous-ghosts}
\end{center}
The full running lab notebook (\texttt{notes.md}), the companion
papers, and the provenance record live in the program's hub
repository,
\url{https://github.com/embeddetech/ghosts-program}.
This paper's results are reproducible from the following scripts,
each standalone (NumPy/SciPy), running in seconds to minutes on a
laptop, and printing its controls inline:
\begin{itemize}\setlength\itemsep{1pt}\raggedright
\item \texttt{qft\_bridge\_stage2.py} --- ladder onset, growth laws,
  channel multiplicity
\item \texttt{qft\_bridge\_stage3.py} --- factorization, dimension
  ladder
\item \texttt{qft\_bridge\_e4\_2d.py} --- the dimensional dichotomy
\item \texttt{qft\_bridge\_decomp.py} --- tadpole/channel
  decomposition
\item \texttt{foliation\_m3.py} --- the frame competition
\item \texttt{s1\_no\_signaling.py},
  \texttt{s2\_infinity\_channel.py},
  \texttt{s3\_foliation\_determinacy.py} --- the signaling,
  completion, and determinacy experiments
\item \texttt{s4a\_books\_balance.py},
  \texttt{s4b\_horn3\_silence.py},
  \texttt{s4c\_dipole\_template.py},
  \texttt{s4d\_cosi\_forecast.py} --- the decision-matrix
  discriminants, the detection template, and the diffuse-background /
  COSI-era detectability forecast
\item \texttt{s5\_state\_divergence.py} --- the state-content
  observable of Claim~\ref{claim:state}
\item \texttt{qft\_bridge\_closeout.py} --- $d=3$ factorization
  Monte Carlo; Garriga--Vilenkin log-growth replication
\item \texttt{s2b\_dirichlet.py} --- the third completion
  (Dirichlet walls) of the infinity-channel taxonomy
\end{itemize}

%=====================================================================
\begin{thebibliography}{99}

\bibitem{indivisible}
J.~Torkelson,
\emph{Indivisible stochastic processes for non-Hermitian and
higher-derivative quantum systems: detailed balance, broken $\PT$
symmetry, and the explosion of ghosts},
companion paper (2026),
\url{https://github.com/embeddetech/p01-indivisible-ghosts}.

\bibitem{Barandes2023}
J.~A.~Barandes,
\emph{The stochastic-quantum correspondence},
Philos.\ Phys.\ \textbf{3}(1), 8 (2025), arXiv:2302.10778.

\bibitem{Barandes2025}
J.~A.~Barandes,
\emph{Quantum systems as indivisible stochastic processes},
arXiv:2507.21192.

\bibitem{PaisUhlenbeck1950}
A.~Pais and G.~E.~Uhlenbeck,
\emph{On field theories with non-localized action},
Phys.\ Rev.\ \textbf{79}, 145 (1950).

\bibitem{Smilga2005}
A.~V.~Smilga,
\emph{Benign vs.\ malicious ghosts in higher-derivative theories},
Nucl.\ Phys.\ B \textbf{706}, 598 (2005).

\bibitem{CHT2003}
S.~M.~Carroll, M.~Hoffman, and M.~Trodden,
\emph{Can the dark energy equation-of-state parameter $w$ be less
than $-1$?},
Phys.\ Rev.\ D \textbf{68}, 023509 (2003).

\bibitem{CJM2004}
J.~M.~Cline, S.~Jeon, and G.~D.~Moore,
\emph{The phantom menaced: constraints on low-energy effective
ghosts},
Phys.\ Rev.\ D \textbf{70}, 043543 (2004), hep-ph/0311312.

\bibitem{GSTZ2021}
C.~Gross, A.~Strumia, D.~Teresi, and M.~Zirilli,
\emph{Is negative kinetic energy metastable?},
Phys.\ Rev.\ D \textbf{103}, 115025 (2021), arXiv:2007.05541.

\bibitem{GV2013}
J.~Garriga and A.~Vilenkin,
\emph{Living with ghosts in Lorentz invariant theories},
JCAP \textbf{01} (2013) 036, arXiv:1202.1239.

\bibitem{Kapner2007}
D.~J.~Kapner, T.~S.~Cook, E.~G.~Adelberger, J.~H.~Gundlach,
B.~R.~Heckel, C.~D.~Hoyle, and H.~E.~Swanson,
\emph{Tests of the gravitational inverse-square law below the
dark-energy length scale},
Phys.\ Rev.\ Lett.\ \textbf{98}, 021101 (2007), arXiv:hep-ph/0611184.

\bibitem{GSV2011}
J.~Garriga, B.~Shlaer, and A.~Vilenkin,
\emph{Minkowski vacua can be metastable},
JCAP \textbf{11} (2011) 035, arXiv:1109.3422.

\bibitem{KOV1975}
I.~Yu.~Kobzarev, L.~B.~Okun, and M.~B.~Voloshin,
\emph{Bubbles in metastable vacuum},
Sov.\ J.\ Nucl.\ Phys.\ \textbf{20}, 644 (1975).

\bibitem{ColemanBounce}
S.~Coleman,
\emph{Fate of the false vacuum: semiclassical theory},
Phys.\ Rev.\ D \textbf{15}, 2929 (1977); erratum \textbf{16}, 1248
(1977);
C.~G.~Callan and S.~Coleman,
Phys.\ Rev.\ D \textbf{16}, 1762 (1977).

\bibitem{Dvali2011}
G.~Dvali,
\emph{Safety of Minkowski vacuum},
arXiv:1107.0956.

\bibitem{DDP2012}
M.~Dine, P.~Draper, and C.~Park,
\emph{Deformed bubbles and Lorentz invariance in vacuum decay},
Phys.\ Rev.\ D \textbf{86}, 065033 (2012), arXiv:1206.5880.

\bibitem{GKT2013}
J.~Garriga, S.~Kanno, and T.~Tanaka,
\emph{Rest frame of bubble nucleation},
JCAP \textbf{06} (2013) 034, arXiv:1304.6681.

\bibitem{GhostCondensate}
N.~Arkani-Hamed, H.-C.~Cheng, M.~A.~Luty, and S.~Mukohyama,
\emph{Ghost condensation and a consistent infrared modification of
gravity},
JHEP \textbf{05} (2004) 074, hep-th/0312099.

\bibitem{Horava-ghosts}
S.~Ramazanov, F.~Arroja, M.~Celoria, S.~Matarrese, and L.~Pilo,
\emph{Living with ghosts in Ho\v rava--Lifshitz gravity},
JHEP \textbf{06} (2016) 020, arXiv:1601.05405.

\bibitem{BenderMannheim2008}
C.~M.~Bender and P.~D.~Mannheim,
\emph{No-ghost theorem for the fourth-order derivative
Pais--Uhlenbeck oscillator model},
Phys.\ Rev.\ Lett.\ \textbf{100}, 110402 (2008).

\bibitem{HawkingHertog}
S.~W.~Hawking and T.~Hertog,
\emph{Living with ghosts},
Phys.\ Rev.\ D \textbf{65}, 103515 (2002), hep-th/0107088.

\bibitem{LHFL2014}
Y.-C.~Lee, M.-H.~Hsieh, S.~T.~Flammia, and R.-K.~Lee,
\emph{Local $\PT$ symmetry violates the no-signaling principle},
Phys.\ Rev.\ Lett.\ \textbf{112}, 130404 (2014).

\bibitem{HypersurfaceBD}
D.~D\"urr, S.~Goldstein, K.~M\"unch-Berndl, and N.~Zangh\`\i,
\emph{Hypersurface Bohm--Dirac models},
Phys.\ Rev.\ A \textbf{60}, 2729 (1999), quant-ph/9801070.

\bibitem{DGNSZ2014}
D.~D\"urr, S.~Goldstein, T.~Norsen, W.~Struyve, and N.~Zangh\`\i,
\emph{Can Bohmian mechanics be made relativistic?},
Proc.\ R.\ Soc.\ A \textbf{470}, 20130699 (2014), arXiv:1307.1714.

\bibitem{Feller1952}
W.~Feller,
\emph{On the integro-differential equations of purely discontinuous
Markoff processes},
Trans.\ Am.\ Math.\ Soc.\ \textbf{48}, 488 (1940);
\emph{The parabolic differential equations and the associated
semi-groups of transformations},
Ann.\ Math.\ \textbf{55}, 468 (1952).

\bibitem{DubovskyNEC}
S.~Dubovsky, T.~Gr\'egoire, A.~Nicolis, and R.~Rattazzi,
\emph{Null energy condition and superluminal propagation},
JHEP \textbf{03} (2006) 025, hep-th/0512260.

\bibitem{BHK2025}
J.~A.~Barandes, M.~Hasan, and D.~Kagan,
\emph{The CHSH game, Tsirelson's bound, and causal locality},
arXiv:2512.18105.

\bibitem{JGB2019}
C.~Jones, T.~Guaita, and A.~Bassi,
\emph{Impossibility of extending the Ghirardi--Rimini--Weber model
to relativistic particles},
Phys.\ Rev.\ A \textbf{103}, 042216 (2021), arXiv:1907.02370.

\bibitem{Tumulka2020}
R.~Tumulka,
\emph{A relativistic GRW flash process with interaction},
in \emph{Do Wave Functions Jump?} (Springer, 2020),
arXiv:2002.00482.

\bibitem{Ovdat2017}
O.~Ovdat, J.~Mao, Y.~Jiang, E.~Y.~Andrei, and E.~Akkermans,
\emph{Observing a scale anomaly and a universal quantum phase
transition in graphene},
Nat.\ Commun.\ \textbf{8}, 507 (2017), arXiv:1701.04121.

\bibitem{KinDipole2026}
T.~T.~Takeuchi,
\emph{A general formulation of the kinematic dipole as a functional
of selection and source properties: beyond the Ellis--Baldwin
approximation},
Astrophys.\ J.\ Suppl.\ \textbf{285}, 3 (2026), arXiv:2602.07389.

\bibitem{LatticeFVD2022}
J.~Braden, M.~C.~Johnson, H.~V.~Peiris, A.~Pontzen, S.~Weinfurtner,
\emph{Mass renormalization in lattice simulations of false vacuum
decay},
Phys.\ Rev.\ D \textbf{107}, 083509 (2023), arXiv:2204.11867.

\bibitem{DESIDR2}
DESI Collaboration,
\emph{DESI DR2 results II: measurements of baryon acoustic
oscillations and cosmological constraints},
Phys.\ Rev.\ D \textbf{112}, 083515 (2025), arXiv:2503.14738.

\bibitem{DESIreview}
T.-N.~Li, G.-H.~Du, H.~Wang, Y.-H.~Li, J.-F.~Zhang, X.~Zhang,
\emph{Dark energy in the DESI era: a brief review of evidence,
beyond-$\Lambda$CDM interpretations, and tensions},
Res.\ Astron.\ Astrophys.\ \textbf{26}, 084002 (2026), arXiv:2606.21826.

\bibitem{DESIskeptic}
M.~Cort\^es, A.~R.~Liddle,
\emph{On DESI's DR2 exclusion of $\Lambda$CDM},
Mon.\ Not.\ R.\ Astron.\ Soc.\ Lett.\ \textbf{544}, L121 (2025), arXiv:2504.15336.

\bibitem{Blackwell1958}
D.~Blackwell,
\emph{Another countable Markov process with only instantaneous
states},
Ann.\ Math.\ Statist.\ \textbf{29}, 313 (1958).

\bibitem{ChungBook}
K.~L.~Chung,
\emph{Markov Chains with Stationary Transition Probabilities},
2nd ed.\ (Springer, Berlin, 1967).

\bibitem{Valentini1991}
A.~Valentini,
\emph{Signal-locality, uncertainty, and the subquantum H-theorem, I
and II},
Phys.\ Lett.\ A \textbf{156}, 5 (1991); \textbf{158}, 1 (1991).

\bibitem{Tang2016}
J.-S.~Tang et al.,
\emph{Experimental investigation of the no-signalling principle in
parity--time symmetric theory using an open quantum system},
Nat.\ Photonics \textbf{10}, 642 (2016).

\bibitem{Brody2016}
D.~C.~Brody,
\emph{Consistency of $\PT$-symmetric quantum mechanics},
J.\ Phys.\ A \textbf{49}, 10LT03 (2016).

\bibitem{Ju2019}
C.-Y.~Ju, A.~Miranowicz, G.-Y.~Chen, and F.~Nori,
\emph{Non-Hermitian Hamiltonians and no-go theorems in quantum
information},
Phys.\ Rev.\ A \textbf{100}, 062118 (2019), arXiv:1906.08071.

\bibitem{EGW2014}
T.~Epelbaum, F.~Gelis, and B.~Wu,
\emph{Nonrenormalizability of the classical statistical
approximation},
Phys.\ Rev.\ D \textbf{90}, 065029 (2014), arXiv:1402.0115.

\bibitem{DGBdipole}
A.~Kashlinsky, F.~Atrio-Barandela, and C.~S.~Shrader,
\emph{Probing the dipole of the diffuse gamma-ray background},
Astrophys.\ J.\ Lett.\ \textbf{961}, L1 (2024), arXiv:2401.04564.

\bibitem{Hardy1992}
L.~Hardy,
\emph{Quantum mechanics, local realistic theories, and
Lorentz-invariant realistic theories},
Phys.\ Rev.\ Lett.\ \textbf{68}, 2981 (1992).

\bibitem{KachelriessSerpico}
M.~Kachelrie{\ss} and P.~D.~Serpico,
\emph{The Compton--Getting effect on ultra-high energy cosmic rays
of cosmological origin},
Phys.\ Lett.\ B \textbf{640}, 225 (2006), astro-ph/0605462.

\bibitem{AmblardDrezet}
A.~Amblard and A.~Drezet,
\emph{Arrival-time distributions as a probe of the preferred
foliation in relativistic Bohmian mechanics},
arXiv:2604.17507.

\bibitem{GeorgiiTumulka}
H.-O.~Georgii and R.~Tumulka,
\emph{Global existence of Bell's time-inhomogeneous jump process
for lattice quantum field theory},
Markov Proc.\ Relat.\ Fields \textbf{11}, 1 (2005), math/0312294.

\bibitem{ChebotarevFagnola}
A.~M.~Chebotarev and F.~Fagnola,
\emph{Sufficient conditions for conservativity of minimal quantum
dynamical semigroups},
J.\ Funct.\ Anal.\ \textbf{153}, 382 (1998), funct-an/9711006.

\bibitem{ClinePhantom2024}
J.~M.~Cline et al.,
\emph{Phantom fluid cosmology: impact of a phantom hidden sector on
cosmological observables},
Phys.\ Rev.\ D \textbf{108}, 095033 (2023), arXiv:2308.12989;
\emph{Boosted dark matter from a phantom fluid},
Phys.\ Lett.\ B \textbf{848}, 138377 (2024), arXiv:2308.01333.

\bibitem{Deffayet2504}
C.~Deffayet, A.~Held, S.~Mukohyama, and A.~Vikman,
\emph{Ghostly interactions in $(1{+}1)$ dimensional classical field
theory},
arXiv:2504.11437.

\bibitem{Deffayet2026}
C.~Deffayet, A.~Fathe Jalali, A.~Held, S.~Mukohyama, and A.~Vikman,
\emph{Unitary time evolution and vacuum for a quantum stable ghost},
arXiv:2604.21823.

\bibitem{KAU2017}
K.~Kawabata, Y.~Ashida, and M.~Ueda,
\emph{Information retrieval and criticality in parity-time-symmetric
systems},
Phys.\ Rev.\ Lett.\ \textbf{119}, 190401 (2017).

\bibitem{Konnig2016}
F.~K\"onnig, H.~Nersisyan, Y.~Akrami, L.~Amendola, and
M.~Zumalac\'arregui,
\emph{A spectre is haunting the cosmos: quantum stability of massive
gravity with ghosts},
JHEP \textbf{11} (2016) 118, arXiv:1605.08757.

\bibitem{Babichev2024}
E.~Babichev,
\emph{Cherenkov radiation as ghost instability},
Phys.\ Rev.\ D \textbf{112}, 084007 (2025), arXiv:2412.20093.

\bibitem{Buoninfante2026}
L.~Buoninfante,
\emph{Ghosts versus unstable particles in quantum field theory},
arXiv:2606.18349.

\bibitem{Andreassen2016}
A.~Andreassen, D.~Farhi, W.~Frost, and M.~D.~Schwartz,
\emph{Precision decay rate calculations in quantum field theory},
Phys.\ Rev.\ D \textbf{95}, 085011 (2017), arXiv:1604.06090.

\bibitem{KaplanSundrum}
D.~E.~Kaplan and R.~Sundrum,
\emph{A symmetry for the cosmological constant},
JHEP \textbf{07} (2006) 042, hep-th/0505265.

\bibitem{Tumulka2006}
R.~Tumulka,
\emph{A relativistic version of the Ghirardi--Rimini--Weber model},
J.\ Stat.\ Phys.\ \textbf{125}, 821 (2006), quant-ph/0406094.

\bibitem{Valentini2007}
A.~Valentini,
\emph{Astrophysical and cosmological tests of quantum theory},
J.\ Phys.\ A \textbf{40}, 3285 (2007), hep-th/0610032.

\bibitem{CG1935}
A.~H.~Compton and I.~A.~Getting,
\emph{An apparent effect of galactic rotation on the intensity of
cosmic rays},
Phys.\ Rev.\ \textbf{47}, 817 (1935).

\end{thebibliography}

\end{document}
